\documentclass[11pt,a4paper]{article}

% Compile with pdfLaTeX. No external images or fonts are required.
\usepackage[utf8]{inputenc}
\usepackage[T1]{fontenc}
\usepackage[french]{babel}
\usepackage[margin=1.65cm]{geometry}
\usepackage{lmodern}
\usepackage{titlesec}
\usepackage{enumitem}
\usepackage[hidelinks]{hyperref}
\usepackage{microtype}

% Searchable, selectable PDF text.
\input{glyphtounicode}
\pdfgentounicode=1

\renewcommand{\familydefault}{\sfdefault}
\pagestyle{empty}
\setlength{\parindent}{0pt}
\setlength{\parskip}{0pt}
\setlength{\emergencystretch}{1em}
\urlstyle{same}
\hypersetup{pdftitle={Reda El Ferrasi - AI - FR},pdfauthor={Reda El Ferrasi}}

\titleformat{\section}{\normalsize\bfseries}{}{0pt}{}[\titlerule]
\titlespacing*{\section}{0pt}{9pt}{5pt}
\setlist[itemize]{leftmargin=1.2em,label=\textbullet,
  topsep=3pt,itemsep=1.5pt,parsep=0pt,partopsep=0pt}

% One line for a role, project or education entry; dates stay on the right.
\newcommand{\entry}[2]{\noindent #1\hfill #2\par}
\newcommand{\jobgap}{\vspace{5pt}}

\begin{document}


\begin{center}
  {\LARGE\bfseries Reda El Ferrasi}\\[3pt]
  {\large Ingénieur IA junior --- agents et intégration LLM}\\[5pt]
  {\small
    \href{tel:+212607420430}{+212 607420430}
    \enspace\textbar{}\enspace
    \href{mailto:redaelferrasi@gmail.com}{redaelferrasi@gmail.com}
    \enspace\textbar{}\enspace
    \href{https://www.linkedin.com/in/reda-elferrasi/}{\nolinkurl{linkedin.com/in/reda-elferrasi/}}
    \enspace\textbar{}\enspace
    \href{https://github.com/redaelfer}{\nolinkurl{github.com/redaelfer}}
  }
\end{center}
\vspace{-5pt}


\section*{PROFIL}
Ingénieur d’État en Data Science et IoT, avec une expérience en orchestration d’agents OpenClaw et Hermes Agent et en déploiement conteneurisé. Projets d’intégration de LLM avec Spring AI et Ollama.

\section*{COMPÉTENCES TECHNIQUES}
\textbf{Programmation :} Python, Java.\par
\textbf{IA :} Spring AI, Ollama.\par
\textbf{Méthodes :} Intégration de LLM, orchestration d'agents, appel d'outils.\par
\textbf{Outils :} Docker, Git.\par
\textbf{Langues :} Arabe (langue maternelle), anglais (B2), français (B2).

\section*{EXPÉRIENCE PROFESSIONNELLE}

\entry{\textbf{Ingénieur IA} \textbar{} Matious Digital -- Projet de fin d'études (PFE)}{janv. 2026 -- sept. 2026}
\begin{itemize}
  \item Développé et testé une plateforme avec \textbf{20 à 40 agents IA simultanés}.
  \item Intégré \textbf{OpenClaw et Hermes Agent} via WebSocket et API REST.
  \item Automatisé la génération de prompts, le choix des modèles, le routage et la reprise sur erreur.
  \item Déployé des agents conteneurisés sur Hetzner Cloud.
\end{itemize}

\jobgap
\entry{\textbf{Stagiaire Data Engineer} \textbar{} ONCF}{juin 2025 -- août 2025}
\begin{itemize}
  \item Automatisé l'ETL SNMP/Syslog : latence réduite de \textbf{23\,\%} lors des tests.
  \item Centralisé les données et créé des tableaux de bord avec \textbf{plus de 10 KPI}.
  \item Réduit de \textbf{35\,\%} le délai de détection des incidents dans les tests NOC.
\end{itemize}

\jobgap
\entry{\textbf{Stagiaire développeur Full Stack} \textbar{} S4U}{juil. 2024 -- août 2024}
\begin{itemize}
  \item Développé un SaaS de tickets React : temps de traitement réduit de \textbf{30\,\%} en test.
  \item Ajouté des tests Jest avec \textbf{72\,\% de couverture} sur les fonctionnalités clés.
\end{itemize}

\section*{PROJETS}

\textbf{AI-Agents} \textbar{} Assistant IA local \textbar{} \href{https://github.com/redaelfer/AI-Agents}{\nolinkurl{github.com/redaelfer/AI-Agents}}\par
\begin{itemize}
  \item Relié un chatbot Streamlit à \textbf{Llama 3.2/Ollama} : appels d'outils Python puis synthèse des résultats.
  \item Limité les extraits de résumé à \textbf{1\,500 caractères} par document et ajouté le suivi \textbf{CPU/RAM} avec psutil.
\end{itemize}

\jobgap
\textbf{TalentAI} \textbar{} Recrutement (React/Spring Boot) \textbar{} \href{https://github.com/redaelfer/TalentAI}{\nolinkurl{github.com/redaelfer/TalentAI}}\par
\begin{itemize}
  \item Conçu les prompts \textbf{Spring AI/Ollama} pour extraire les compétences, évaluer la compatibilité CV/offre et générer des questions.
  \item Testé le scoring et l'extraction avec \textbf{deux tests JUnit/Mockito}, sur des réponses du LLM simulées.
\end{itemize}

\jobgap
\textbf{RetinaScan} \textbar{} IA médicale \textbar{} \href{https://github.com/redaelfer/RetinaScan}{\nolinkurl{github.com/redaelfer/RetinaScan}}\par
\begin{itemize}
  \item Adapté \textbf{MobileNetV2} sous \textbf{TensorFlow} : base gelée, augmentation d'images et \textbf{dropout (0,5)}.
  \item Suivi l'exactitude entraînement/validation sur \textbf{10 époques} ; exposé les prédictions via une API Flask.
\end{itemize}

\section*{FORMATION}
\entry{\textbf{ENSIAS} -- Diplôme d'ingénieur d'État -- Data Science \& IoT}{2023 -- 2026}
\vspace{3pt}
\entry{\textbf{CPGE} -- TSI}{2021 -- 2023}

\section*{CERTIFICATIONS \& CENTRES D'INTÉRÊT}
\textbf{CS50x} \textbar{} Harvard University\par
\vspace{3pt}
\textbf{AWS Knowledge: Cloud Essentials} -- Badge de formation \textbar{} AWS\par

\vspace{5pt}
\textbf{Intérêts :} Reprogrammation ECU, mécanique/réparation auto-moto, modding de jeux, histoire/mythologie.

\end{document}
